\documentclass[../../main.tex]{subfiles}
\graphicspath{{../../image/}{../../figure/}}
\begin{document}

\section{注记}\label{section:ch09-注记}

如历史导言所指出,Riemann 问题最初由 Riemann [1] 在等熵气体动力学方程组\eqref{align:ch07-等熵Euler}的框架下提出并求解。

关于\cref{section:ch09-自相似解与riemann问题}所勾勒的双曲守恒律一般自相似解结构,以及真正非线性方程组中此类解必为有界变差特殊函数的证明,读者可参阅 Dafermos [28] 与 Elling、Roberts [2]。

由 Riemann 构想的激波与稀疏波曲线方法逐步发展,用以求解等熵或绝热气体动力学方程组的特殊 Riemann 问题,描述波相互作用与激波管实验。这项早期研究综述于 Courant、Friedrichs [1]。该工作的提炼促成 Lax [2] 对严格双曲、特征族真正非线性或线性退化的守恒律方程组(Theorem 9.4.1)以弱波求解 Riemann 问题。绝热(非等熵)气体动力学方程组 Riemann 问题解的详细论述见 Smoller [3]、Serre [11]、Godlewski、Raviart [2]、Holden、Risebro [5] 的教材,特别地见 Chang、Hsiao [3] 的专著。讨论大初值问题的早期文献包括 Smoller [1,2]、Smith [1] 与 Sever [1,2,3]。特别地,Sever [3] 警示:即使对真正非线性方程组,标准激波容许准则也可能不足以保证大初值 Riemann 问题解的唯一性。

处理非真正非线性方程组需额外努力。遵循 Oleinik $E$ 条件的范式,一般标量守恒律 Riemann 问题解的形式由 Gelfand [1] 通过一个例子描述。随后,Wendroff [1] 在 $\sigma_{uu}$ 可变号时求解了方程组\eqref{align:ch07-p系统}与\eqref{align:ch07-热弹性}的 Riemann 问题。\eqref{align:ch07-标量守恒律}与\eqref{align:ch07-p系统}中借助 $f$ 与 $\sigma$ 凸包或凹包的解构造,见 Dafermos [2] 与 Leibovich [1]。上述结果显然早在 1960 年代已为中国的研究预见,只是后来才传至国际科学界,例如 Chang、Hsiao [1,2] 及 Hsiao、Zhang [1]。Chang、Hsiao [3] 的专著给出详细论述与大量文献。特别地,Chen、Young [2] 研究了具一般状态方程的气体非等熵流方程组 Riemann 问题解中是否会出现真空状态。对特殊方程组的研究,推动了 Tai-Ping Liu [1] 对任意分段真正非线性、严格双曲守恒律方程组 Riemann 问题的求解。Iguchi、LeFloch [1] 对与分段真正非线性特征族相伴的波曲线结构作了更精细分析,从而求解了更一般的、不必分段真正非线性的严格双曲守恒律方程组 Riemann 问题。另见 Ancona、Marson [4]。波扇曲线可能仅 Lipschitz 连续这一观察归功于 Bianchini [6]。

对流通量函数仅 Lipschitz 连续的方程组,Riemann 问题由 Correia、LeFloch、Thanh [1] 求解。对流通量函数有跳跃间断的方程组,见 LeFloch、Thanh [2,4]。另见 Holden、Risebro [4], Chalons、Raviart、Seguin [1], Colombo、Marcellini [1] 及 Andreianov [1]。平衡律方程组的 Riemann 问题在 Hong、Temple [1] 中讨论,亦见 Goatin、LeFloch [4]。

关于 $\delta$ 激波的早期文献,见 Keyfitz、Kranzer [2,3,4]。详细论述见 Sever [12] 的备忘录,以及 Keyfitz [4] 的"回顾与展望"。在很大程度上, $\delta$ 激波理论是与无压气体模型研究相伴发展的。在此联系下,见 Li、Zhang、Yang [1], Feimin Huang [2], Sever [10,11] 及\cref{section:ch11-注记}所引文献。此外,Gui-Qiang Chen、Hailiang Liu [1,2] 证明:当压强对密度的响应松弛到零时,等熵或非等熵气体动力学方程组 Riemann 问题的解,在极限下退化为含 $\delta$ 激波的无压气体动力学方程组的解。$\delta$ 激波亦由 Ercole [1], Joseph [1], Hayes、LeFloch [1], Tan [1], Tan、Zhang、Zheng [1], Sheng、Zhang [1], Hanchun Yang [1], Li、Yang [1], Brenier [3], Danilov、Shelkovich [1,2], Danilov、Mitrovich [2], Panov、Shelkovich [1], Shelkovich [1,2,3], Keyfitz、Tsikkou [1], Tsikkou [3] 及 Joseph、Sahoo [2] 讨论。

熵率容许准则由 Dafermos [3] 提出。关于来自动理学的动机,见 Ferziger、Kaper [1,\S5.5] 与 Kohler [1]。消粘性方法提供了进一步动机;见 Bethuel、Despres、Smets [1]。真正非线性标量守恒律熵率准则刻画容许 $L^{\infty}$ 解的证明,见 Blaser、Rivière [1]。Theorems 9.7.2 与 9.7.5 取自 Dafermos [32],Theorem 9.7.6 归功于 Hsiao [1]。Theorem 9.7.3 的详细替代证明见 Dafermos [18]。最后,Theorem 9.7.4 取自 Dafermos [29]。

熵率准则的效力还在变更类型、刻画相变的方程组上得到检验(Hattori [1,2,3,4,5,6,7], Pence [2])。另见 Aavatsmark [1], Sever [7] 及 Krejčí、Straskraba [1]。

二维 Euler 方程组熵率准则的状况在 Chiodaroli、Kreml [1] 中讨论。

由熵不等式给出 Riemann 问题解的另一种刻画,归功于 Heibig、Serre [1]。关于熵极小化的其他结果,见 Tadmor [1], Kröner、LeFloch、Thanh [1], Belletini、Bertini、Mariani、Novega [1], 及 Bonaschi、Carillo、DiFrancesco、Peletier [1]。

把双曲守恒律方程组的自相似解作为带时变人工粘性的耗散方程组自相似解极限来研究,由 Kalasnikov [1]、Tupciev [1,2] 与 Dafermos [4] 各自独立开创。该方法最初用于 Dafermos [4,5] 及 Dafermos、DiPerna [1] 求解两个守恒律严格双曲方程组的 Riemann 问题。特别地,此处对方程组\eqref{align:ch09-9-8-33}的处理改编自 Dafermos、DiPerna [1]。

Theorem 9.8.1 归功于 Tzavaras [3],它最早完整构造了一般严格双曲守恒律方程组 Riemann 问题的解,不要求分段真正非线性。同一般水平的平行构造,用标准消粘性方法(见第 XV 章),后由 Bianchini、Bressan [5] 与 Bianchini [6] 给出。事实上,后一文献启发了此处\cref{section:ch09-粘性波扇}中发展的粘性波扇法构造波曲线。用粘性波扇法构造波曲线的另一种途径,见 Joseph、LeFloch [4]。

关于用粘性波扇法构造特定双曲方程组 Riemann 问题解,读者可参阅下列文献:严格双曲方程组见 Yong Jung Kim [1], Slemrod、Tzavaras [1,3], Tzavaras [2,5], Andreianov [2];耦合方程组见 Boutin、Coquel、LeFloch;含 $\delta$ 激波的非严格双曲方程组见 Ercole [1], Keyfitz、Kranzer [3,4], Joseph [1], Tan [1], Tan、Zhang、Zheng [1], Sheng、Zhang [1], Li、Yang [1];混合型方程组见 Slemrod [4], Fan [1,2];等熵气体动力学球对称解见 Slemrod [7]。

用标准消粘性方法显式构造等熵或非等熵气体流方程组 Riemann 问题解,见 Lin、Yang [1], Huang、Jian、Wang [1], Huang、Wang、Yang [1,2], Huang、Li、Wang [1], Huang、Li [1], Huang、Wang、Wang、Yang [1], Zhang、Pan、Wang、Tan [1], Zhang、Pan、Tan [1]。

关于 Riemann 问题结构稳定解(即使存在强激波和/或欠压激波)能否用粘性波扇逼近的问题,见 Schecter [4,5,6], Lin、Schecter [1], Schecter、Szmolyan [1], Weishi Liu [1], Xiao-Biao Lin [1,2,3,4,5,6]。相关数值计算见 Schecter、Marchesin、Plohr [3]。

改造守恒律方程组标准扰动以保持解自相似性这一想法有众多推广。在此联系下,见 Li、Zhang、Yang [1](二维无压气体动力学方程组), Slemrod、Tzavaras [2] 与 Tzavaras [1](气体动力学 Broadwell 方程组), Joseph、LeFloch [1,2,3,4,5] 及 LeFloch、Rohde [1](半平面或四分之一平面上的 Riemann 问题,一般粘性矩阵,或以时变粘性毛细作用或松弛代替粘性);Jiequan Li、Peng Zhang [1](模拟燃烧的方程组\eqref{align:ch07-燃烧模型},常反应速率 $\delta$ 换为 $\delta/t$)。

非严格双曲方程组 Riemann 问题理论的现状远未定论。解的存在性与容许性都提出棘手问题,因为波扇可含形形色色奇异波,如过压激波、欠压激波、$\delta$ 激波与振荡波。致力于包罗万象的理论是徒劳的;应聚焦连续物理中出现的特定方程组,尤其是弹性力学与多相流。此类方程组的分类以及容许解存在唯一性方面已有进展;见 Glimm [2], Azevedo、Marchesin [1], Azevedo、Marchesin、Plohr、Zumbrun [1,2], Freistühler [3], Isaacson、Marchesin、Plohr [1], Isaacson、Marchesin、Plohr、Temple [1,2], Isaacson、Temple [2], Schaeffer、Shearer [2], Schecter、Marchesin、Plohr [1,2], M. Shearer [4,5], Shearer、Schaeffer、Marchesin、Paes-Leme [1], Schecter、Shearer [1], Schulze、Shearer [1], Tang、Ting [1], Zhu、Ting [1], Čanić [1,2], Čanić、Peters [1], Ercole [1], Keyfitz、Kranzer [1,2,3], Tan [1], Tan、Zhang、Zheng [1], Schecter [1,2,3], Chalons、Coquel [1], Hanche-Olsen、Holden、Risebro [1], Azevedo、de Souza、Furtado、Marchesin [1], Azevedo、Eschenazi、Marchesin、Palmeira [1], Lambert、Marchesin [2], Rodrigues-Bermúdez、Marchesin [1], Shapiro、Marchesin、Schecter [1], Silva、Marchesin [1]。Ostrov [3] 给出一个富有启发性的例子:在该例中,用消粘性方法通过改变粘性系数相对大小,得到 Riemann 问题的不同解。

读者可从 Tong Zhang、Yuxi Zheng [2] 对控制 Chapman-Jouguet 燃烧的平衡律方程组的处理,略窥连续物理中方程组 Riemann 问题解所遇波构型之丰。微分几何中也有有趣的 Riemann 问题;见 Bascar、Prasad [1]。

用于模拟相变的混合型方程组 Riemann 问题解,除经典激波与稀疏波外还含相边界。如\cref{section:ch08-注记}所注,相边界的容许性由动力学关系决定。此类 Riemann 问题的解见 Fan [1,2,5,6], Frid、Liu [1], Hattori [1,2,3,4,7,8], Holden [1], Hsiao [2], Hsiao、DeMottoni [1], Keyfitz [1], LeFloch、Thanh [1], Mercier、Piccoli [1,2], Pence [2], M. Shearer [1,3], Shearer、Yang [1], Slemrod [5], Corli、Fan [1], Fan、Lin [1]。富有教益的讨论与更多文献见 LeFloch [5] 的专著。

Riemann 问题也已对非守恒形式的拟线性双曲方程组\eqref{align:ch09-9-8-18}提出,并通过\cref{section:ch08-非守恒激波}所述途径拼接稀疏波与激波求解(见 LeFloch、Liu [1], LeFloch、Tzavaras [1], Colombo、Monti [1]),或用消粘性方法求解(Bianchini、Bressan [5], Bianchini [6])。R. Young [10] 给出了那种形式的方程组,其 Riemann 问题容许含稀疏波而无激波的多个解。

四分之一平面上 Riemann 问题用消粘性方法的求解,见 Bianchini、Spinolo [1]。同样可用\cref{section:ch09-粘性波扇}所述粘性波扇构造,见 Christoforou、Spinolo [1,2,3]。应用见 Andrianov、Warnecke [1,2], LeFloch、Thanh [2], LeFloch、Shearer [1], Joseph、Sahoo [1]。

所谓\emph{广义 Riemann 问题},即初值在跳跃间断两侧光滑(而非常数),见 Ta-tsien Li、De-yin Kong [2], Ta-tsien Li、Libin Wang [2,4,5,6], Mentrelli、Ruggeri [1], Ben Artzi [1], Ben-Artzi、Li [1]。全面论述见 Ben-Artzi、Falcovitz [1] 的专著。

对真正非线性方程组,波相互作用研究与 Theorem 9.9.1 的原始证明归功于 Glimm [1]。此处给出的推导取自 Yong [1]。对非真正非线性方程组,波相互作用估计最初由 Tai-Ping Liu [15] 得到,他首先认识到入射角的关键作用。近期详细而严格的论述见 Iguchi、LeFloch [1], Tai-Ping Liu、Tong Yang [7]。(非线性)弹性弦中波相互作用的讨论见 R. Young [5],相对论理论中见 Groah、Smoller、Temple [1]。等熵气体动力学中产生出奇异特征解的复合波相互作用,见 R. Young [9]。非等熵理想气体方程组可能波相互作用的完整分类见 Chang、Hsiao [3],以及 Chen、Endres、Jenssen [1]。实际波相互作用的描述:等熵弹性方程组见 Greenberg [1,2],标量燃烧模型见 Liu、Zhang [1],带摩擦阻尼等熵气流 Euler 方程组见 Luo、Yang [1]。

\cref{section:ch09-弱解的崩坏}所呈弱解崩坏例子取自 Jenssen [1]。更多例子由 Baiti、Jenssen [3], R. Young [5,6], Young、Szeliga [1], Jenssen、Young [1] 构造。特别地,文中证明:当解的值域上特征速度族并不一致分离时,即使初值全变差任意小,解也可能在有限时间内崩坏。关于解上确界或全变差快速放大乃至爆炸的早期工作,见 Jeffrey [2], R. Young [2], Joly、Métivier、Rauch [2]。迄今为止,在具物理兴趣的方程组解中尚未探测到爆炸,这升起希望:此类方程组的特殊结构可能抵消不稳定源。但另见 Tsikkou [1],以及 Bressan、Chen、Zhang、Zhu [1]。

\end{document}
